\documentclass[aspectratio=169]{beamer}
\usetheme{Madrid}
\usecolortheme{default}

% Packages
\usepackage{graphicx}
\usepackage{booktabs}
\usepackage{hyperref}
\usepackage{amsmath}
\usepackage{listings}
\usepackage{xcolor}

\definecolor{codegreen}{rgb}{0,0.6,0}
\definecolor{codegray}{rgb}{0.5,0.5,0.5}
\definecolor{codepurple}{rgb}{0.58,0,0.82}
\definecolor{backcolour}{rgb}{0.95,0.95,0.92}

\lstdefinestyle{mystyle}{
    backgroundcolor=\color{backcolour},   
    commentstyle=\color{codegreen},
    keywordstyle=\color{magenta},
    numberstyle=\tiny\color{codegray},
    stringstyle=\color{codepurple},
    basicstyle=\ttfamily\footnotesize,
    breakatwhitespace=false,         
    breaklines=true,                 
    captionpos=b,                    
    keepspaces=true,                 
    numbers=left,                    
    numbersep=5pt,                  
    showspaces=false,                
    showstringspaces=false,
    showtabs=false,                  
    tabsize=2
}
\lstset{style=mystyle}

% Title Page Information
\title[RL for Financial Reasoning]{Reinforcement Learning Fine-Tuning of a Small Language Model for Faithful Financial Reasoning over Financial Reports}
\author[Maheshwari et al.]{
    Vedant Maheshwari - CB.SC.U4AIE23346 \and \\
    Prisha Gupta - CB.SC.U4AIE23346 \and \\
    Naksh Singh - CB.SC.U4AIE23361 \and \\
    Akhil - CB.SC.U4AIE23317
}
\institute{Academic Project Presentation}
\date{\today}

\begin{document}

\begin{frame}
    \titlepage
\end{frame}

\begin{frame}{Outline}
    \tableofcontents
\end{frame}

\section{Motivation \& Problem Statement}
\begin{frame}{Motivation: Current Problem}
    Large Language Models are increasingly used for investment analysis, financial report understanding, and ratio analysis.
    
    \vspace{0.3cm}
    \textbf{However, they frequently:}
    \begin{itemize}
        \item Hallucinate financial values
        \item Perform incorrect calculations
        \item Use wrong figures from tables
        \item Ignore numerical consistency
        \item Produce logically inconsistent reasoning
    \end{itemize}
\end{frame}

\begin{frame}{Why SFT is Insufficient}
    Supervised Fine-Tuning (SFT) primarily teaches the model format and style, but \textbf{it does not optimize}:
    \begin{itemize}
        \item Numerical accuracy
        \item Reasoning correctness
        \item Equation correctness
        \item Consistency across intermediate steps
    \end{itemize}
    \vspace{0.3cm}
    \begin{block}{Problem Statement}
        Existing SLMs fine-tuned on financial QA datasets often generate numerically inconsistent answers despite fluent text. This project proposes an RL-based framework to directly optimize reasoning and numerical consistency.
    \end{block}
\end{frame}

\section{Objectives \& Outcomes}
\begin{frame}{Objective \& Expected Outcomes}
    \textbf{Objective:} Develop a reinforcement learning framework that improves numerical reasoning, arithmetic accuracy, financial reasoning, and evidence grounding—while preserving language fluency.
    
    \vspace{0.3cm}
    \textbf{Expected Outcomes (without increasing model size):}
    \begin{itemize}
        \item Improve Exact Match / Accuracy (Primary Metric)
        \item Improve F1 Score
        \item Reduce hallucinations
        \item Improve reasoning correctness
    \end{itemize}
\end{frame}

\section{Datasets \& Models}
\begin{frame}{Datasets}
    \textbf{Currently Completed:}
    \begin{itemize}
        \item \textbf{FinQA:} A large-scale dataset with 2.8k financial reports for 8k Q\&A pairs requiring multi-step numerical reasoning over structured and unstructured evidence.
    \end{itemize}
    
    \vspace{0.3cm}
    \textbf{Next Phases:}
    \begin{itemize}
        \item \textbf{ConvFinQA:} Multi-turn conversations for sequential reasoning.
        \item \textbf{TAT-QA:} Tabular And Textual dataset containing text, tables, and arithmetic reasoning.
    \end{itemize}
\end{frame}

\begin{frame}{Models \& Reward Function}
    \textbf{Models Evaluated:}
    \begin{itemize}
        \item \textbf{Qwen2.5-1.5B-Instruct} (Current Focus)
        \item SmolLM2-1.7B
        \item Qwen2.5-0.5B
    \end{itemize}
    
    \vspace{0.3cm}
    \textbf{Composite Reward Function Design:}
    \begin{itemize}
        \item \textbf{Answer Correctness:} Does the final output exactly match the gold float/string?
        \item \textbf{Equation Correctness:} Is the generated mathematical DSL valid and executable?
        \item \textbf{Evidence Grounding:} Did the model extract the correct numbers?
        \item \textbf{Reasoning Completeness:} Does the model provide a step-by-step breakdown?
    \end{itemize}
\end{frame}

\section{Implementation Challenges}
\begin{frame}{Challenges Discovered During RL Training}
    During the initial evaluation, Exact Match Accuracy plateaued at 18\%. An in-depth analysis revealed critical environment flaws:
    \begin{enumerate}
        \item \textbf{DSL Operator Mismatch:} The evaluation engine failed to parse raw operators from the training data (e.g., \texttt{add1-1}, \texttt{divide1-2}) and variables (\texttt{const\_100}), returning 0 reward for correct reasoning.
        \item \textbf{Binary Classification (Yes/No):} The \texttt{greater()} operator evaluated to strings (\texttt{yes/no}), crashing the evaluator which expected float outputs.
        \item \textbf{Reward Saturation:} The \texttt{grounding\_reward} saturated at 1.0 with 0 variance, providing no learning signal and diluting the execution reward budget.
    \end{enumerate}
\end{frame}

\begin{frame}{Solutions Implemented}
    To fix the reward signal, we completely overhauled the RL environment:
    \begin{itemize}
        \item \textbf{Execution Engine Rewrite:} Added support for constant handling (\texttt{const\_X}), percentage stripping, and dynamic operator normalization.
        \item \textbf{Yes/No Support:} Modified the system prompt to explicitly teach the model binary comparison formatting and updated the evaluator to handle string comparisons gracefully.
        \item \textbf{Reward Rebalancing:} Stripped the saturated \texttt{grounding\_reward} to concentrate the KL-divergence budget strictly on execution accuracy and equation correctness.
    \end{itemize}
\end{frame}

\section{Results \& Evaluation}
\begin{frame}{Results: Accuracy Improvement Across Runs}
    \begin{table}[]
        \centering
        \begin{tabular}{lcc}
            \toprule
            \textbf{Training Phase} & \textbf{Overall Accuracy} & \textbf{Yes/No Accuracy} \\
            \midrule
            Epoch 1 (SFT + RL, pre-fixes) & $\sim 18.00\%$ & $0.00\%$ (Crashes) \\
            Epoch 3 (RL End-to-End, post-fixes) & \textbf{46.64\%} & \textbf{55.00\%} (11/20) \\
            \bottomrule
        \end{tabular}
        \caption{Evaluation on FinQA validation set using Qwen 1.5B.}
    \end{table}
    
    \vspace{0.3cm}
    \textbf{Key Takeaway:} Fixing the reward environment and binary operator logic caused a massive \textbf{+28\% jump} in overall execution accuracy.
\end{frame}

\begin{frame}{Comparison with Original FinQA Baselines}
    \textbf{How does our 46.64\% compare to the original FinQA paper?}
    \begin{itemize}
        \item \textbf{Original FinQANet (RoBERTa-Large):} Achieved $\sim 61.24\%$ \textbf{but} relied on a dedicated \textit{Retriever} model to filter out 99\% of the text beforehand.
        \item \textbf{Our RL Model (Qwen 1.5B):} Achieves 46.64\% \textbf{End-to-End} without a retriever.
    \end{itemize}
    \vspace{0.3cm}
    \begin{block}{Conclusion}
        For a small 1.5B model performing both retrieval and reasoning simultaneously from raw context, crossing 45\% accuracy is a highly competitive academic result that proves RL improves numerical reasoning.
    \end{block}
\end{frame}

\begin{frame}
    \begin{center}
        \Huge \textbf{Thank You!} \\
        \vspace{1cm}
        \normalsize Questions?
    \end{center}
\end{frame}

\end{document}
